% ----------------------------------------------------------
% Introdução
% ----------------------------------------------------------
\chapter{Introdução}

A presença das Tecnologias Digitais de Informação e Comunicação (TDIC) vem causando mudanças bastante profundas nas maneiras como a cultura é produzida, acessada e preservada hoje em dia. Estudos do Nic \cite{nic-cultura-tecnologias-no-brasil} mostram, por exemplo, que, apesar do crescimento significativo no uso de dispositivos conectados para atividades culturais (como assistir vídeos, ouvir música ou acessar conteúdos educativos), ainda existe uma desigualdade bem marcante na distribuição dessas ferramentas digitais que valorizam o patrimônio local. Essa lacuna fica ainda mais visível quando falamos de cidades menores ou médias, onde as manifestações culturais e as memórias coletivas muitas vezes continuam praticamente invisíveis, limitando o seu reconhecimento e também dificultando que sejam preservadas. \cite{prolam-cultura-viva-2025} \cite{soster2021tecnologias}

Pensando nesse cenário, o objetivo deste trabalho é analisar e demonstrar como aplicativos móveis podem realmente se transformar em ferramentas úteis para democratizar o acesso à cultura regional. Essa motivação tem muito a ver com o fato de que hoje aproximadamente 80\% das casas brasileiras têm acesso à internet, sendo que 62\% usam exclusivamente dispositivos móveis pra isso \cite{mcom-2023-internet}. Então, considerando essa popularidade dos smartphones, aparece aí uma oportunidade real de superar barreiras sociais e geográficas, facilitando bastante a promoção, registro e atualização constante das informações sobre patrimônios culturais e históricos locais.

A importância dessa discussão fica ainda mais clara quando observamos que muitas publicações acadêmicas brasileiras \cite{prolam-cultura-viva-2025} destacam que políticas públicas e iniciativas locais de preservação cultural muitas vezes não possuem ferramentas tecnológicas eficientes, de baixo custo e que facilitem a participação popular e a inclusão digital. Isso acaba dificultando bastante a divulgação das memórias coletivas. Nesse contexto, os aplicativos móveis podem ajudar muito ao oferecer funcionalidades práticas, como geolocalização, upload de imagens, filtros de acessibilidade e categorização regional, permitindo que a própria comunidade participe ativamente na valorização cultural da sua região.

Para conseguir abordar esse assunto da melhor maneira possível, este trabalho está organizado em quatro etapas principais. Primeiro, será feita uma fundamentação teórica sobre a relação entre cultura e tecnologia, destacando algumas experiências brasileiras recentes que utilizam inovação digital para promoção cultural. Em seguida, os objetivos específicos da pesquisa serão detalhados, explicando também a escolha da região e das tecnologias utilizadas. Depois, será descrito o desenvolvimento do aplicativo, mostrando claramente sua arquitetura e as funcionalidades principais, explicando como essas escolhas ajudam a resolver as necessidades identificadas. Por fim, será feita uma análise crítica dos resultados obtidos, discutindo o potencial impacto dessa solução na ampliação do acesso à cultura e no fortalecimento da identidade local.

Dessa forma, o trabalho não busca só contribuir teoricamente para o debate sobre cultura digital e democratização do patrimônio, mas também apresentar uma proposta prática baseada em dados concretos e experiências regionais, que pode ajudar bastante outras iniciativas futuras de inclusão cultural mediadas pelas tecnologias digitais no Brasil.


\section{Contextualização}

A cultura, que pode ser entendida como uma expressão viva da identidade, memória e das práticas sociais, hoje em dia está diretamente relacionada às tecnologias digitais. No Brasil atual, fica cada vez mais claro que vivemos numa sociedade conectada, onde o acesso à informação e às redes digitais transforma a forma como as pessoas e os grupos registram, consomem e compartilham suas experiências culturais \cite{nic-cultura-tecnologias-no-brasil}. A popularização dos smartphones e das redes móveis aumenta ainda mais esse fenômeno, tornando muito comum o uso de aplicativos no cotidiano para comunicação, diversão e também para acessar serviços relacionados à cultura.

Por outro lado, essa conectividade não está distribuída igualmente pelo país. Enquanto as grandes cidades têm um acesso mais fácil às plataformas digitais culturais, muitas cidades menores ou regiões periféricas ainda enfrentam dificuldades para registrar, divulgar e valorizar seus patrimônios culturais. Segundo o relatório "Cultura e Tecnologias no Brasil" \cite{nic-cultura-tecnologias-no-brasil}, a maior parte das iniciativas de digitalização cultural está concentrada principalmente nas capitais e nas cidades turísticas, deixando de fora muitas manifestações e bens culturais das regiões menos visíveis, que acabam ficando marginalizadas ou até mesmo esquecidas.

Ao mesmo tempo, algumas políticas públicas nacionais, como o programa Cultura Viva \cite{prolam-cultura-viva-2025}, já mostram que há um reconhecimento institucional sobre a importância de conectar o universo digital à gestão, preservação e divulgação da cultura. Essas iniciativas buscam incentivar redes de Pontos de Cultura e estimular o uso de tecnologias digitais por grupos e comunidades locais, mas ainda enfrentam desafios importantes, como infraestrutura, capacitação das pessoas envolvidas e continuidade das ações.

Nesse cenário, então, surge uma demanda cada vez maior por soluções tecnológicas inclusivas e de fácil acesso, que realmente considerem as realidades específicas dos diferentes territórios brasileiros. Os aplicativos móveis, que oferecem recursos como geolocalização, registro de imagens, opções de acessibilidade e filtros regionais, destacam-se como ferramentas estratégicas capazes de envolver a população local, democratizar o acesso à cultura e fortalecer a identidade regional. \cite{soster2021tecnologias} \cite{oliveira-cultura-tecnologias-ufse}

Portanto, discutir o uso de aplicativos móveis para mapear e divulgar pontos culturais não é só uma questão tecnológica, mas está diretamente ligada a debates mais amplos, como o direito à memória, a valorização do patrimônio, a participação da população e a inclusão digital. Tudo isso faz parte de um esforço coletivo para transformar o acesso à cultura em uma política pública mais ampla e efetiva, com alcance nacional, usando soluções digitais adaptadas às realidades diversas do Brasil.


\section{Justificativa}

Este estudo se justifica pela necessidade de dar mais atenção e valor ao patrimônio cultural local, principalmente em regiões que costumam receber menos investimentos e visibilidade. Diversas manifestações culturais acabam sendo esquecidas ou pouco divulgadas, o que enfraquece a memória e a identidade das comunidades. De acordo com o IBGE e estudos da área, o Brasil possui políticas de preservação cultural desde 1937\cite{ibge-internet-2023}. No entanto, essas ações ainda são aplicadas de forma desigual e não chegam com eficiência a muitas cidades menores.

Além disso, dados do Ministério das Comunicações mostram que cerca de 80\% das residências no país têm acesso à internet. Desse total, 62\% das pessoas usam somente dispositivos móveis, como celulares, para se conectar\cite{mcom-2023-internet}. Esse dado reforça a importância dos aplicativos móveis como ferramentas estratégicas para divulgar informações culturais, já que grande parte da população utiliza o celular como principal meio de acesso.

Outro ponto importante é a falta de plataformas digitais específicas para mapear e divulgar a cultura local. São poucas as iniciativas que oferecem, em um só lugar, dados como localização, imagens, descrições e informações de acessibilidade dos pontos culturais. Essa ausência dificulta o conhecimento e o reconhecimento da cultura local pela própria população, além de limitar ações educativas e turísticas que poderiam ser desenvolvidas nesses espaços.

Diante desse cenário, a proposta de criação de um aplicativo para cadastro de pontos culturais aparece como uma solução relevante e viável. O aplicativo vai permitir que o usuário acesse, de forma simples e organizada, conteúdos como localização no mapa, fotos, descrições e informações sobre acessibilidade dos espaços culturais do estado.

Por fim, o uso de uma tecnologia acessível e inclusiva pode ampliar o acesso à cultura, ajudar a preservar as identidades locais e contribuir para a inclusão digital. Dessa forma, o aplicativo poderá ser uma ferramenta importante para promover o reconhecimento, o pertencimento e a valorização dos espaços culturais pelas próprias comunidades.

\section{Objetivos}
\subsection{Objetivo Geral}

Analisar como o uso de um aplicativo de cadastro de pontos culturais pode colaborar com o fortalecimento e a promoção do desenvolvimento cultural no estado.

\subsection{Objetivos Específicos}
\begin{lista}
  \item Desenvolver um aplicativo destinado ao cadastro de pontos culturais;
  \item Realizar o mapeamento de locais de interesse cultural, incluindo imagens, descrições e localização geográfica;
  \item Inserir informações relativas à acessibilidade com o intuito de ampliar o acesso à cultura;
  \item Promover a divulgação de pontos turísticos e históricos para a população em geral.
\end{lista}



\section*{Resumo}

O capítulo introdutório apresentou o contexto da transformação digital dentro da área cultural, mostrando como as tecnologias móveis podem ser ferramentas importantes para apoiar e fortalecer as manifestações culturais locais. Também foram definidos o objetivo geral e os objetivos específicos da pesquisa, com ênfase no desenvolvimento de um aplicativo voltado para o mapeamento e a divulgação de pontos culturais. A proposta busca promover mais inclusão, acessibilidade e valorização do patrimônio cultural das regiões envolvidas.

Os próximos capítulos estão organizados da seguinte forma:

\begin{lista}
  \item \textbf{Fundamentação Teórica:} Neste capítulo são apresentados todos os conceitos teóricos utilizados no desenvolvimento da solução proposta no presente trabalho;
  \item \textbf{Trabalhos Relacionados} Apresentando projetos que possuem objetivos similares a este que esta sendo desenvolvido;
  \item \textbf{Sobre o Aplicativo} Capítulo dedicado à apresentação do protótipo da solução implementada, detalhando funcionalidades, utilização e demais detalhes da implementação;
  \item \textbf{Conclusão:} Neste capítulo é feita a conclusão do trabalho, dados seus objetivos propostos, e são listados os trabalhos futuros para melhorar e/ou expandir a utilização da solução.
\end{lista}